%% language-de-sigplan-test — German proceedings paper (acmart `language=german` on sigplan).
%% Proceedings formats print the abstract heading (\abstractname → "Zusammenfassung") that journal formats omit, and every format heads the bibliography with \refname ("Literatur"); both come from babel.
%% Also covers \keywordsname ("Schlagwörter"), \proofname ("Beweis") and \acksname ("Danksagungen").
%% Twin of language-de-sigplan-test.typ.
\documentclass[sigplan,language=german]{acmart}

\setcopyright{acmlicensed}
\copyrightyear{2018}
\acmYear{2018}
\acmMonth{6}
\acmDOI{XXXXXXX.XXXXXXX}
\acmISBN{978-1-4503-XXXX-X/2018/06}
\acmConference[Conference'17]{Proceedings of ACM Conference}{June 2018}{Washington, DC, USA}
\acmBooktitle{Proceedings of ACM Conference (Conference'17)}

\begin{document}

\title{Eine Notiz über Berechnungskomplexität}

\author{Hans Müller}
\email{mueller@example.de}
\affiliation{%
  \institution{Forschungsinstitut}
  \city{Berlin}
  \country{Germany}}

\begin{abstract}
  Eine kurze Notiz auf Deutsch, um die lokalisierten Überschriften eines Tagungsbeitrags zu prüfen.
\end{abstract}

\keywords{Komplexität, Algorithmen, Berechnung}

\maketitle

\section{Einleitung}
Dieses Dokument prüft die lokalisierten Zeichenketten eines Tagungsbeitrags: die Überschrift der Zusammenfassung, die Schlagwörter, den Beweis, die Danksagung und die Überschrift des Literaturverzeichnisses~\cite{Cohen07}.

\begin{proof}
Ein kurzer Beweis, der mit einem QED-Quadrat endet.
\end{proof}

\begin{acks}
Wir danken den Gutachtern.
\end{acks}

\bibliographystyle{ACM-Reference-Format}
\bibliography{sample-base}

\end{document}
